In this section, we describe the steps needed to compile your mutator and mutatee programs and to run
them. First we give you an overview of the major steps and then we explain each one in detail.

\subsection{Creating a mutator program}

The first step in using Dyninst is to create a mutator program. The mutator program specifies the
mutatee (either by naming an executable to start or by supplying a process ID for an existing
process). In addition, your mutator will include the calls to the API library to modify the mutatee.
For the rest of this section, we assume that the mutator is the sample program given in Appendix A.

The following \code{CMakeLists.txt} shows how create your mutator program with Dyninst.

\begin{lstlisting}
  cmake_minimum_required(VERSION 3.15)
  project(retee LANGUAGES CXX)
  
  find_package(Dyninst REQUIRED)
  
  add_executable(retee retee.cpp)
  target_link_libraries(retee PUBLIC Dyninst::dyninstAPI)
\end{lstlisting}

\subsection{Setting Up the Application Program (mutatee)}

On most platforms, any additional code that your mutator might need to call in the mutatee (for example
files containing instrumentation functions that were too complex to write directly using the API) can
be put into a dynamically loaded shared library, which your mutator program can load into the mutatee
at runtime using the loadLibrary member function of \code{BPatch\_process}.

To locate the runtime library that Dyninst needs to load into your program, an additional environment
variable must be set. The variable \code{DYNINSTAPI\_RT\_LIB} should be set to the full pathname of the
run time instrumentation library, which should be:


NOTE: \code{DYNINST\_LIB} should be set to the location where Dyninst library files were installed

\begin{lstlisting}
\$(DYNINST_LIB)/libdyninstAPI_RT.so (UNIX)

\%DYNINST_LIB/libdyninstAPI_RT.dll (Windows)
\end{lstlisting}

\subsection{Running the Mutator}

At this point, you should be ready to run your application program with your mutator. For example,
to start the sample program shown in Appendix A:

\code{retee foo <pid>}
